\documentclass[10pt,a4paper]{amsart}
\usepackage[margin=19mm]{geometry}
\usepackage{amsmath,amssymb,mathtools}
\usepackage[scr=rsfs,cal=euler]{mathalfa}
\usepackage{xfrac}
\usepackage[unicode,hidelinks,bookmarks=false]{hyperref}
\newcommand{\ie}{i.e.\ }
\newcommand{\cf}{cf.\ }
\newcommand{\resp}{respectively\ }
\newcommand{\Subsection}{\subsection}
\providecommand{\nobreakdash}{\nobreak\hbox{-}}
\setlength{\emergencystretch}{2em}
\multlinegap0pt
\usepackage{amssymb}
\usepackage{xcolor}
\newtheorem{proposition}{Proposition}
\newcommand{\NN}{\mathbf{N}}
\def\dd{\mathrm{d}}
\title{A fixed-point theorem for face maps,\\ or deletion-tolerant random finite sets}
\author{Tom Hutchcroft, Nicolas Monod, Omer Tamuz}
\date{Source: arXiv:2505.21484v2 (CC BY 4.0)}
\begin{document}
\maketitle

\noindent\emph{Source and licence.} A fixed-point theorem for face maps,\\ or deletion-tolerant random finite sets. Version arXiv:2505.21484v2 (CC BY 4.0), distributed by arXiv under the Creative Commons Attribution 4.0 International licence, http://creativecommons.org/licenses/by/4.0/, as recorded in the arXiv licence metadata for this exact version and shown on the abstract page https://arxiv.org/abs/2505.21484v2. Full licence notice: \texttt{licenses/arxiv-2505.21484-CC-BY-4.0.txt}.\par\smallskip

\noindent\textit{Editorial scope.} Counted unit: the article's proposition prop:ergodic (source0.tex lines 610--613). The article's complete original proof of it is delivered with it (source0.tex lines 615--673, one \texttt{proof} environment of 59 lines). No mathematics is written, repaired or re-derived: the statement and the proof are the article's own lines, in the article's own notation, and no retained line is substituted or re-flowed.  No further source-local context is needed: every cross-reference of every retained line resolves inside the two delivered spans.  Nothing was omitted from any retained span, so the package carries no omission record.  No retained line contains a citation, so the wrapper carries no bibliography and no bibliography entry.  No retained line contains a figure environment, an image-inclusion command or drawing code, so no image is delivered and the package declares no asset dependency.  Editorial formatting only, in addition: an AMS \texttt{amsart} preamble that restates the article's own declarations and macros used by the retained lines, namely the environments \texttt{proposition} on separate counters, together with the standard packages those lines need.  The retained environments print in this document as Proposition 1 in order of appearance, whereas the article prints the same items with the numbering of its own full document.  The complete native source of the article as distributed in the arXiv e-print for this exact version is bundled unchanged as \texttt{sources/178873/source0.tex}.
\begin{proposition}
    \label{prop:ergodic}
    Every invariant ergodic probability measure on $X$ is IID.
\end{proposition}

\begin{proof}
    Let $\mu$ be invariant and ergodic. Let $\pi_1 \colon \Omega \to X$ be the projection $\pi_1(\omega)=\omega_1$, and let $\eta = \pi_1\mu$ be the law of $\omega_1$. That is, for measurable $Y \subseteq X$, $\eta(Y) = \mu(\{\omega\,:\, \omega_1 \in Y\})$. 
    
    We will show that if $Y_1,Y_2,\ldots$ are subsets of $X$ then 
    \begin{align}
        \label{eq:product}
        \mu\left(\prod_i Y_i\right) = \prod_i \eta(Y_i).
    \end{align}

    Fix any measurable $Y \subseteq X$. For $i \in \NN$, let
    \begin{align*}
      F_i = \{\omega \,:\, \omega_i \in Y\}.
    \end{align*}
    Since $F_i = s_1^{1-i}F_1$ and since $\mu$ is $s_1$-invariant, $\mu(F_i) = \mu(F_1) = \eta(Y)$.

    Fix any measurable $W \subseteq X^{M-1}$, and let
    \begin{align*}
      G = \{\omega\,:\, (\omega_1,\ldots,\omega_{M-1}) \in W\}.
    \end{align*}
    We claim that $\mu(F_M \cap G) = \mu(F_M) \mu(G)$. This will then imply~\eqref{eq:product} by a standard argument. Indeed, by the sigma-additivity of probability measures it suffices to show that if $Y_1,\ldots,Y_M$ are subsets of $X$, and if $A_i = \{\omega\,:\,\omega_i \in Y_i\}$, then $\mu(A_i\cap \cdots \cap A_M) = \eta(Y_1)\cdots\eta(Y_M)$. To this end, it suffices to show that $\mu(A_i\cap \cdots \cap A_M) = \mu(A_i\cap \cdots \cap A_{M-1})\eta(Y_M)$, as the proof then follows by induction.

Turning to the claim, let $f_i \in L^2(\Omega,\mu)$ be given by
    \begin{align*}
      f_i(\omega) = 1_{F_i}(\omega)-\mu(F_i) = 1_{F_i}(\omega)-\mu(F_1).
    \end{align*}
    Note that $\int f_i\,\dd\mu = 0$ and $\| f_i \|^2 =  \mu(F_1)(1-\mu(F_1)) \leq 1$. For $i\in \NN$, let $U_i$ be the isometric operator on $L^2(\Omega,\mu)$ given by $U_ig = g \circ s_i$. It follows that
    \begin{align*}
      U_i f_n =
      \begin{cases}
        f_n&\text{if } n < i\\
        f_{n+1}&\text{otherwise.}
      \end{cases}
    \end{align*}

    By the von Neumann ergodic theorem the limit
    \begin{align*}
      f = \lim_N \frac{1}{N}\sum_{k=0}^{N-1} U_M^k f_M
    \end{align*}
    exists and is equal to the projection of $f_M$ on the space of $U_M$-invariant elements of $L^2(\Omega,\mu)$. Since $U_M^k f_M = f_{M+k}$, we have that
    \begin{align*}
      f = \lim_N \frac{1}{N}\sum_{k=0}^{N-1} f_{M+k} = \lim_N \frac{1}{N}\sum_{k=0}^{N-1} f_{k},
    \end{align*}
    where the second equality follows from the fact that
    \begin{align*}
      \left\| \frac{1}{N}\sum_{k=0}^{N-1} f_{M+k} - \frac{1}{N}\sum_{k=0}^{N-1} f_{k} \right\| \leq \frac{2M}{N}.
    \end{align*}
    By the same argument, it holds for any $i \in \NN$ that
    \begin{align*}
      f = \lim_N \frac{1}{N}\sum_{k=0}^{N-1} f_{i+k} = \lim_N \frac{1}{N}\sum_{k=0}^{N-1} U_i^kf_{i+k}.
    \end{align*}
    Thus $f$ is $U_i$-invariant for all $i$, and it follows from ergodicity that $f$ is constant. Moreover, $\int f_{i}\,\dd\mu = 0$, and so $\int f\,\dd\mu = 0$. Thus $f=0$.

    Let $g(\omega) = 1_G(\omega)-\mu(G)$, and note that $U_Mg=g$, by the definition of $G$. Now,
    \begin{multline*}
      (g,f_M) = \int (1_G(\omega)-\mu(G))(1_{F_M}(\omega) - \mu(F_M))\,\dd\mu(\omega)\\
      = \mu(G \cap F_M) - \mu(G)\mu(F_M).
    \end{multline*}
    Since $U_M g = g$ and $U_M$ preserves inner products, it holds that $(g,U_M^kf_M) = (g,f_M)$. Thus also $(g,f)=\mu(G \cap F_M) - \mu(G)\mu(F_M)$. Since $f=0$, it follows that $\mu(G \cap F_M) = \mu(G)\mu(F_M)$ as claimed.
\end{proof}

\end{document}
